% TOP_PROBLEM: {"id": 1255, "title": "Infinitely Many Giuga Numbers", "category_id": 1, "category": {"id": 1, "name": "number_theory", "display_name": "Number Theory", "description": "Properties of integers, prime numbers, Diophantine equations.", "slug": "number-theory", "order_index": 1, "created_at": "2026-07-31T15:26:25.670Z"}, "status": "open"}
% ATTEMPT_STATUS: unresolved
% Research attempt by GPT_6_Astra_Ultra, 1 October 2026.
% Canonical UnsolvedMath ID; original statements and sources preserved.
% FIRST_POSTED: 2026-10-01
% Imported from: unsolvedmath_additional_500.tex; UnsolvedMath ID 1255
% !TeX program = lualatex
% Compile twice: lualatex 1255.tex
% Self-contained: no external bibliography, images, or \input files.
% Numeric section labels and visible numbers are original UnsolvedMath IDs.
\documentclass[11pt,a4paper]{article}
\usepackage[margin=24mm,headheight=14pt]{geometry}
\usepackage{fontspec}
\setmainfont{Latin Modern Roman}
\setsansfont{Latin Modern Sans}
\setmonofont{Latin Modern Mono}
\usepackage{amsmath,amssymb,mathtools}
\usepackage{unicode-math}
\setmathfont{Latin Modern Math}
\usepackage{microtype}
\usepackage{enumitem,xurl,needspace}
\usepackage[dvipsnames]{xcolor}
\usepackage{hyperref}
\hypersetup{unicode=true,colorlinks=true,linkcolor=MidnightBlue,urlcolor=MidnightBlue,
 pdftitle={UnsolvedMath Problem 1255},pdfauthor={GPT 6 Astra Ultra},bookmarksnumbered=true}
\usepackage{bookmark,tocloft,titlesec,fancyhdr}
\setlength{\cftsecnumwidth}{4.2em}
\cftsetpnumwidth{2.5em}
\cftsetrmarg{4em}
\titleformat{\section}{\Large\bfseries\raggedright}{\thesection}{0.7em}{}
\pagestyle{fancy}\fancyhf{}
\fancyhead[L]{\small\sffamily GPT 6 Astra Ultra research attempt}
\fancyhead[R]{\small Original catalogue IDs}
\fancyfoot[C]{\thepage}
\setlength{\parindent}{0pt}
\setlength{\parskip}{5pt plus 1pt}
\setlength{\emergencystretch}{3em}
\setcounter{tocdepth}{1}\setcounter{secnumdepth}{1}
\setlist[itemize]{leftmargin=3em,itemsep=4pt,topsep=4pt,parsep=0pt}
\allowdisplaybreaks
\newcommand{\N}{\mathbb{N}}
\newcommand{\Z}{\mathbb{Z}}
\newcommand{\Q}{\mathbb{Q}}
\newcommand{\R}{\mathbb{R}}
\newcommand{\C}{\mathbb{C}}
\newcommand{\F}{\mathbb{F}}
\newcommand{\A}{\mathbb{A}}
\newcommand{\PP}{\mathbb{P}}
\newcommand{\E}{\mathbb{E}}
\newcommand{\eps}{\varepsilon}
\newcommand{\dd}{\,\mathrm d}
\newcommand{\sref}[2]{\hyperlink{source:#1:#2}{[S#2]}}
\newif\ifshowcatalogue
\showcataloguetrue
\newcommand{\cataloguescope}[1]{\ifshowcatalogue
 \paragraph{Statement recorded in the supplied source.}
 #1\fi}
\begin{document}
% UnsolvedMath ID 1255; source catalogue code NT-048

\setcounter{section}{1254}

\section{Infinitely Many Giuga Numbers}\label{1255}

{\small\sffamily UnsolvedMath ID 1255 · Number Theory}

\subsection{Definitions and mathematical statement}

\paragraph{Shared notation.}Unless a section says otherwise, $\N=\{1,2,\ldots\}$,
$\Z$ is the ring of integers, $\Q,\R,\C$ are the rational, real, and complex
fields, $[N]=\{1,\ldots,\lfloor N\rfloor\}$, and $\log$ is the natural
logarithm. For real functions of a parameter tending to infinity,
$f\ll g$ and $f=O(g)$ mean $|f|\leq Cg$ eventually for some fixed $C$;
$f\gg g$ reverses this comparison; $f=o(g)$ means $f/g\to0$;
$f\sim g$ means $f/g\to1$; and $f\asymp g$ means both order bounds.
A subscript on $O$ or $\ll$ specifies allowed constant dependence.
The expression $n^{o(1)}$ in an upper bound means growth smaller than
$n^\epsilon$ for each fixed $\epsilon>0$. Local definitions govern any
exception, finite-field convention, or use of the same symbol for another
object. An asymptotic claim is not automatically an effective algorithm.



A Giuga number is a composite integer $n$ such that $p\mid(n/p-1)$ for every prime divisor $p$ of $n$. This condition forces $n$ to be squarefree. It is stronger than merely satisfying a congruence for one of its prime factors. The target is an unbounded family of composite integers meeting all these divisibility conditions. \sref{1255}{1} \sref{1255}{2}

\paragraph{Statement recorded in the supplied source.}
Are there infinitely many Giuga numbers? \sref{1255}{1}

\paragraph{Residual task reported by the archive.}

The embedded review (2026-08-17) records: Fix the intended Giuga convention, then prove an infinite family or a finiteness result for that precisely stated class. \sref{1255}{1}

\subsection{Short English statement}

Are there infinitely many composite integers for which dividing by any prime factor leaves a number one greater than a multiple of that same prime?

\subsection{Research attempt}

\paragraph{Convention and source check.}
This is the infinitude question for composite integers satisfying
$p\mid n/p-1$ at every prime divisor, without an additional
Carmichael condition. Luca--Pomerance--Shparlinski [S2] was inspected
in its author-hosted primary PDF on 1 October 2026. The following
reciprocal identity, small-support classification, and extension
obstruction are direct deductions for exactly this convention.

\paragraph{Squarefreeness and an equivalent reciprocal condition.}
If $p^2\mid n$, then $n/p-1\equiv-1\pmod p$, impossible. Thus
$n$ is squarefree. Set $P=\{p:p\mid n\}$ and
\[
 H=\sum_{p\in P}\frac np-1.
\]
Modulo a prime $r\mid n$, every term in the sum except $n/r$ vanishes,
so $H\equiv n/r-1\pmod r$. Since the prime divisors are distinct,
all the defining congruences hold if and only if $n\mid H$.
Equivalently,
\[
                   \sum_{p\mid n}\frac1p-\frac1n\in\mathbb Z.
\]
For composite squarefree $n$, this quantity is positive. Indeed there
are at least two terms $n/p$ greater than one in its numerator.
Hence it is an integer at least one for every Giuga number.
The identity is an equivalence only together with squarefreeness;
that hypothesis was established before applying the Chinese remainder
argument.

\paragraph{Complete classification with at most three prime factors.}
There are no one-prime-factor examples because a composite prime power
is not squarefree. If $n=pq$ with $p<q$, the condition at $q$ would
give $q\mid p-1$, impossible. Now take $n=pqr$ with $p<q<r$.
The positive integer in the preceding identity is
\[
                  T=\frac1p+\frac1q+\frac1r-\frac1{pqr}.
\]
If $p\ge3$, then $T<1/3+1/5+1/7<1$, contradiction. Thus $p=2$.
If $q\ge5$, then $T<1/2+1/5+1/7<1$, again impossible. Hence
$q=3$, and
\[
                         T=\frac56+\frac5{6r}.
\]
For $r\ge5$ this is at most one, with equality exactly when $r=5$.
Thus thirty is the unique Giuga number with at most three distinct
prime factors. Its congruences are explicit:
$2\mid14$, $3\mid9$, and $5\mid5$.

\paragraph{Appending one new prime can never preserve the property.}
Suppose $n$ is Giuga and $q\nmid n$ is a new prime. If $nq$ were
Giuga too, then for every $p\mid n$ we would have both
$n/p\equiv1\pmod p$ and $q(n/p)\equiv1\pmod p$. Hence
$q\equiv1\pmod p$ for every $p\mid n$, so $q\equiv1\pmod n$.
In particular $q\ge n+1$. But the defining condition at the new
prime is $q\mid n-1$, impossible for a positive $n-1<q$.
Therefore no Giuga number can be extended to another just by
multiplying by a single new prime.

This rules out a natural Euclid-style recursion. It does not rule out
adding several primes together or replacing part of the support.
For an extension by a squarefree coprime integer $Q$, the old-prime
conditions similarly force $Q\equiv1\pmod n$, but each new prime
condition depends on all the other new factors. That coupled system
is not covered by the one-prime contradiction.

\paragraph{The unit-fraction search and its uniformity problem.}
For a fixed number $k$ of prime factors, the integer reciprocal sum
places useful bounds on the smallest primes. More generally after
fixing $p_1<\cdots<p_j$, every remaining reciprocal is at most
$1/p_{j+1}$, so an upper estimate for the remaining sum can prune
some branches. However, infinitude quantifies over unbounded $k$.
A finite enumeration at fixed $k$ or fixed largest prime cannot prove
finiteness of the entire class unless one also bounds those parameters.

The correction $1/n$ also cannot simply be discarded. It is exactly
what turns the modular conditions into an integer, and changing it to
zero produces a different Egyptian-fraction equation. An approximate
unit-fraction sum near an integer does not produce divisibility at
every prime factor of the product.

\paragraph{Executed exact check.}
The script factors every $4\le n\le200000$, rejects prime powers
and all nonsquarefree inputs, and checks every required divisibility.
It obtains exactly
\[
                         30,\quad858,\quad1722,\quad66198.
\]
These are finite examples and not a new construction of infinitely many.
Code and output are \texttt{checks\_second.py} and
\texttt{results\_second.json} in
\path{_Local/Erdos_problems/UnsolvedMath/attempt_audit_2026_10_01/number_theory/}.
The exact search uses the catalogue's congruences, not a Bernoulli
or prime-characterization convention.

\paragraph{Remaining gap and outcome.}
The attempt proves the reciprocal equivalence, the three-prime-factor
classification, and a complete obstruction to single-prime extensions.
No unbounded construction survives these constraints, but no theorem
bounds all supports either. \textbf{Outcome: UNRESOLVED.} Eliminating
one plausible recursive construction does not prove finiteness.

\subsection{Sources}

{\small

\begin{itemize}

\item[\hypertarget{source:1255:1}{S1}] ULAM.AI, \emph{UnsolvedMath}, supplied \texttt{problems.json}, record 1255 (NT-048), statement and background. The embedded literature-review date is 2026-08-17; that is the archive’s review date, not a fresh certification. Dataset landing page: \url{https://huggingface.co/datasets/ulamai/UnsolvedMath}.

\item[\hypertarget{source:1255:2}{S2}] Florian Luca, Carl Pomerance, and Igor E. Shparlinski, On Giuga Numbers, International Journal of Modern Mathematics 4 (2009), 13-18. \url{https://gauss.dartmouth.edu/~carlp/giugafinal.pdf}. Bibliographic information imported from the supplied record.

\item[\hypertarget{source:1255:3}{S3}] OEIS A007850, Giuga numbers. \url{https://oeis.org/A007850}. Bibliographic information imported from the supplied record.


\item[E1] Florian Luca, Carl Pomerance, and Igor E. Shparlinski,
\emph{On Giuga Numbers}, author-hosted primary paper.
\url{https://gauss.dartmouth.edu/~carlp/giugafinal.pdf}.
Inspected 1 October 2026. The congruence convention is retained; no
literature bound is needed for the direct arguments above.

\end{itemize}

}

\label{end:1255}
\end{document}
